\chapter{Attributes}
\label{ch:attributes}

\begin{quote}
\textit{``Metadata is a love note to the future.''} --- Jason Scott
\end{quote}

Attributes provide a way to attach metadata to code elements---functions, structs, enums, and variables. The compiler uses this metadata to alter behavior, generate code, or produce warnings. This chapter covers all implemented attributes with practical examples.

\section{Understanding Attributes}
\label{sec:understanding-attributes}

In systems programming, we often need to tell the compiler things that can't be expressed in the code itself. Should this function be inlined? Is this struct packed for hardware access? Is this function safe to call from an interrupt handler?

Attributes provide these hints and directives. Unlike comments, attributes are part of the language---the compiler parses them, validates them, and acts on them.

\begin{comingfromc}
\textbf{GCC/Clang Attributes:} C has \texttt{\_\_attribute\_\_((packed))}, \texttt{\_\_attribute\_\_((always\_inline))}, etc. Novus uses cleaner \texttt{@packed} and \texttt{@inline} syntax. Same ideas, better ergonomics.
\end{comingfromc}

\section{Attribute Syntax}
\label{sec:attribute-syntax}

Novus supports two equivalent attribute syntaxes. Choose whichever feels more natural for your codebase.

\subsection{At-Sign Syntax (Recommended)}

The primary syntax uses \texttt{@} followed by the attribute name:

\begin{lstlisting}
@test
fn test_addition() {
    expect_eq(2 + 2, 4, "basic addition")
}

@inline
fn fast_add(a: i32, b: i32) -> i32 {
    return a + b
}
\end{lstlisting}

This syntax is concise and visually distinct from code.

\subsection{Bracket Syntax}

An alternative bracket syntax is also supported:

\begin{lstlisting}
#[test]
fn test_multiplication() {
    expect_eq(3 * 7, 21, "multiplication")
}

#[derive(Eq, Hash)]
struct Point {
    x: i32,
    y: i32,
}
\end{lstlisting}

The bracket syntax is useful when attributes have complex arguments.

\subsection{Attribute Arguments}

Attributes can accept arguments---both named and positional:

\begin{lstlisting}
// Named arguments
@library(name = "example.library", version = 1, revision = 0)
pub struct ExampleLibrary {
    counter: u32,
}

// Positional arguments
@test("Verify sorting performance")
fn test_sort_perf() {
    // ...
}

// Flag arguments
@test(skip = "Not yet implemented")
fn test_future() {
    // Skipped during test runs
}

// Mixed arguments
@libfunc(offset = -30)  // LVO offset -30 bytes
pub fn my_library_function() -> i32 {
    return 42
}
\end{lstlisting}

\subsection{Multiple Attributes}

Multiple attributes can be stacked on a single item:

\begin{lstlisting}
@test
@inline
pub fn test_fast_path() {
    // This function is both a test case and should be inlined
}

@interrupt
@export
pub fn vblank_handler() {
    // Interrupt handler that must be preserved in the binary
}
\end{lstlisting}

The order of attributes generally doesn't matter, though some combinations are invalid (the compiler will warn you).

\section{Testing Attributes}
\label{sec:testing-attributes}

Novus has first-class support for testing through attributes.

\subsection{@test}

Marks a function as a test case. The \texttt{novus test} command discovers and compiles these into a test runner.

\begin{lstlisting}
@test
fn test_basic() {
    expect_eq(1 + 1, 2, "one plus one")
}

@test("Descriptive test name")
fn test_with_description() {
    expect_eq(true, true, "truth is true")
}
\end{lstlisting}

\textbf{Arguments:}
\begin{description}
    \item[positional string] Optional test description
    \item[\texttt{skip = "reason"}] Skip this test with an explanation
\end{description}

\begin{lstlisting}
@test(skip = "Hardware not available in CI")
fn test_blitter_operations() {
    // This test requires real Amiga hardware
}

@test(skip)
fn test_known_broken() {
    // Skipped without explanation (not recommended)
}
\end{lstlisting}

\subsection{@bench / @benchmark}

Marks a function for performance testing. Both names are equivalent:

\begin{lstlisting}
@bench
fn bench_sort() {
    var data = generate_random_data(1000)
    sort(&var data)
}

@benchmark
fn bench_hash() {
    var map = HashMap::new()
    for i in 0..1000 {
        map.insert(i, i * 2)
    }
}
\end{lstlisting}

Benchmark functions are included when you run \texttt{novus test -b}. The runner executes each benchmark multiple times and reports timing statistics.

\section{Optimization Attributes}
\label{sec:optimization-attributes}

These attributes guide the compiler's optimization decisions.

\subsection{@inline}

Hints that the compiler should inline this function at call sites:

\begin{lstlisting}
@inline
fn min(a: i32, b: i32) -> i32 {
    if a < b { return a } else { return b }
}

@inline
fn abs(x: i32) -> i32 {
    if x < 0 { return -x } else { return x }
}
\end{lstlisting}

Small, frequently-called functions benefit most from inlining. Benefits include:
\begin{itemize}
    \item Elimination of function call overhead
    \item Better register allocation across the inlined code
    \item Opportunities for constant folding when arguments are known
\end{itemize}

The compiler may ignore this hint for:
\begin{itemize}
    \item Recursive functions
    \item Very large functions
    \item Functions with complex control flow
\end{itemize}

\subsection{@noinline}

Prevents inlining even when the compiler might otherwise choose to inline:

\begin{lstlisting}
@noinline
fn error_handler() {
    // Keep error handling out of hot paths
    // to improve instruction cache locality
    log_error("Something went wrong")
    cleanup_resources()
}

@noinline
fn cold_path_check(value: i32) -> bool {
    // Rarely-taken branch - don't pollute hot code
    return validate_complex_condition(value)
}
\end{lstlisting}

Use \texttt{@noinline} for:
\begin{itemize}
    \item Error handling code (rarely executed)
    \item Large functions called from tight loops
    \item Functions where you need predictable stack usage
\end{itemize}

\section{Layout Attributes}
\label{sec:layout-attributes}

\subsection{@packed}

Removes padding between struct fields:

\begin{lstlisting}
// Without @packed: 8 bytes (a=1, padding=3, b=4)
struct Normal {
    a: u8,
    b: u32,
}

// With @packed: 5 bytes (a=1, b=4, no padding)
@packed
struct Packed {
    a: u8,
    b: u32,
}
\end{lstlisting}

\begin{comingfromc}
\textbf{C Equivalent:} \texttt{\#pragma pack(1)} or \texttt{\_\_attribute\_\_((packed))}. Same effect, but Novus makes it explicit per-struct.
\end{comingfromc}

\textbf{Warning:} Packed structs may have unaligned fields, which can cause:
\begin{itemize}
    \item Slower memory access (multiple bus cycles)
    \item Address errors on 68000/68010 (crashes!)
    \item Correct but slow access on 68020+
\end{itemize}

Use when:
\begin{itemize}
    \item Matching external data formats (file formats, network protocols)
    \item Mapping hardware registers
    \item Memory-constrained situations where every byte matters
\end{itemize}

\section{Trait Derivation}
\label{sec:derive-attributes}

\subsection{@derive / \#[derive]}

Automatically generates trait implementations for structs:

\begin{lstlisting}
#[derive(Eq)]
struct Point {
    x: i32,
    y: i32,
}

#[derive(Eq, Hash)]
struct UserId {
    id: u32,
}
\end{lstlisting}

\textbf{Supported traits:}
\begin{description}
    \item[\texttt{Eq}] Generates \texttt{eq(\&self, \&Self) -> bool} comparing all fields
    \item[\texttt{Hash}] Generates \texttt{hash(\&self) -> u32} combining field hashes
\end{description}

You can derive multiple traits at once:

\begin{lstlisting}
#[derive(Eq, Hash)]
struct CacheKey {
    module: u32,
    offset: u32,
}

// Now CacheKey can be used as HashMap key
var cache = HashMap::<CacheKey, Data>::new()
\end{lstlisting}

\section{Method Chaining}
\label{sec:chain-attribute}

\subsection{@chain}

Enables fluent method chaining by making a method return \texttt{\&var Self}:

\begin{lstlisting}
struct WindowBuilder {
    width: i32,
    height: i32,
    title: *u8,
    flags: u32,
}

impl WindowBuilder {
    pub fn new() -> WindowBuilder {
        return WindowBuilder {
            width: 640,
            height: 480,
            title: "Window",
            flags: 0,
        }
    }

    @chain
    pub fn width(&var self, w: i32) {
        self.width = w
    }

    @chain
    pub fn height(&var self, h: i32) {
        self.height = h
    }

    @chain
    pub fn title(&var self, t: *u8) {
        self.title = t
    }

    pub fn build(&self) -> Window {
        // Create actual window...
    }
}
\end{lstlisting}

Usage:

\begin{lstlisting}
var window = WindowBuilder::new()
    .width(800)
    .height(600)
    .title("My Application")
    .build()
\end{lstlisting}

The \texttt{@chain} attribute automatically adds an implicit \texttt{return self} at the end of the method.

\section{Interoperability Attributes}
\label{sec:interop-attributes}

\subsection{@export}

Marks a function for export with an unmangled name, preventing dead code elimination:

\begin{lstlisting}
@export
pub fn my_callback(data: *u8) -> i32 {
    // Can be called from external C code or AmigaOS
    return process(data)
}

@export
pub fn _UserLibInit() -> i32 {
    // Entry point called by AmigaOS
    return initialize_library()
}
\end{lstlisting}

Use this when:
\begin{itemize}
    \item Creating callback functions for AmigaOS (hooks, interrupt handlers)
    \item Functions that need to be called from assembly or C
    \item Entry points that might appear unused to the optimizer
\end{itemize}

\subsection{@extern\_type}

Marks a struct as an external C type, preventing Novus from emitting its definition:

\begin{lstlisting}
@extern_type
struct timeval {
    tv_sec: i32,
    tv_usec: i32,
}

@extern_type
struct Library {
    node: Node,
    flags: u8,
    // ... matches C header definition
}
\end{lstlisting}

This tells the compiler: ``This type is defined elsewhere (in C headers or runtime); just use it, don't emit a definition.''

\section{Synchronization Attributes}
\label{sec:sync-attributes}

These attributes provide automatic synchronization wrappers.

\subsection{@atomic}

Wraps the function body in a Forbid/Permit critical section:

\begin{lstlisting}
@atomic
pub fn update_shared_counter(counter: *u32) {
    // Forbid() called on entry
    *counter = *counter + 1
    // Permit() called on exit (even if returning early)
}

@atomic
pub fn swap_values(a: *i32, b: *i32) {
    let temp = *a
    *a = *b
    *b = temp
    // No task can switch during this operation
}
\end{lstlisting}

This prevents task switching while the function executes. The compiler ensures Permit() is called on all exit paths.

\textbf{Use for:}
\begin{itemize}
    \item Short, fast operations on shared data
    \item Lock-free data structure updates
    \item Counter increments visible to multiple tasks
\end{itemize}

\textbf{Avoid for:}
\begin{itemize}
    \item Functions that might block (I/O, waiting)
    \item Long-running computations
    \item Anything that allocates memory
\end{itemize}

\subsection{@no\_interrupts}

Wraps the function body in Disable/Enable:

\begin{lstlisting}
@no_interrupts
pub fn poke_hardware_register(addr: *u16, value: u16) {
    // Disable() called on entry - ALL interrupts blocked
    *addr = value
    // Enable() called on exit
}

@no_interrupts
pub fn read_modify_write(reg: *u16, mask: u16, value: u16) {
    let current = *reg
    *reg = (current & ~mask) | (value & mask)
}
\end{lstlisting}

\textbf{Warning:} Use sparingly! This blocks \textit{all} interrupts including:
\begin{itemize}
    \item Audio DMA (causes audio glitches)
    \item Disk I/O (can corrupt data)
    \item Vertical blank (freezes display updates)
    \item Timer interrupts (breaks timing)
\end{itemize}

Only use for timing-critical hardware access that must be atomic.

\begin{comingfromc}
\textbf{AmigaOS:} Disable() disables all hardware interrupts by setting the processor interrupt mask. Enable() restores it. These are Exec functions, not direct hardware access.
\end{comingfromc}

\section{Interrupt Handling Attributes}
\label{sec:interrupt-attributes}

These attributes are essential for writing interrupt-driven code.

\subsection{@interrupt}

Marks a function as an interrupt handler. The compiler generates appropriate prologue/epilogue code:

\begin{lstlisting}
@interrupt
pub fn vblank_handler() {
    // Called during vertical blank interrupt
    // All registers properly saved/restored
    update_frame_counter()
}

@interrupt
pub fn audio_handler() {
    // Called when audio channel needs more data
    feed_next_sample()
}
\end{lstlisting}

The compiler generates:
\begin{itemize}
    \item Register save (MOVEM.L to stack)
    \item Your function body
    \item Register restore (MOVEM.L from stack)
    \item RTE (Return from Exception)
\end{itemize}

\subsection{@interrupt\_safe}

Documents that a function is safe to call from interrupt context:

\begin{lstlisting}
@interrupt_safe
pub fn increment_counter(counter: *u32) {
    // Safe: no allocations, no blocking calls
    *counter = *counter + 1
}

@interrupt_safe
fn calculate_checksum(data: *u8, len: u32) -> u32 {
    // Safe: pure computation, no system calls
    var sum: u32 = 0
    for i in 0..len {
        sum = sum + (data[i] as u32)
    }
    return sum
}
\end{lstlisting}

Functions marked \texttt{@interrupt\_safe} must:
\begin{itemize}
    \item Never allocate memory
    \item Never call blocking functions (Wait, WaitIO, etc.)
    \item Never call Permit() without matching Forbid()
    \item Complete quickly (microseconds, not milliseconds)
\end{itemize}

This is currently documentation-only but helps communicate intent and catch errors during code review.

\section{Diagnostic Attributes}
\label{sec:diagnostic-attributes}

\subsection{@suppress}

Suppresses specific compiler warnings:

\begin{lstlisting}
@suppress("W0002", "Known safe: buffer size validated externally")
pub fn write_unchecked(buffer: *u8, data: *u8, len: i32) {
    // Warning W0002 (unchecked buffer access) suppressed
    // We accept responsibility for bounds checking
}

@suppress("W0001")  // Minimum suppression - just the code
pub fn use_deprecated_api() {
    call_old_function()
}
\end{lstlisting}

\textbf{Arguments:}
\begin{description}
    \item[warning code] The warning ID to suppress (e.g., \texttt{"W0001"})
    \item[justification] Explanation of why the warning is safe to ignore (recommended)
\end{description}

Always prefer fixing warnings over suppressing them. When suppression is necessary, document why.

\section{Program Configuration}
\label{sec:program-attributes}

\subsection{\#[stack\_size]}

Sets the stack size for the compiled program:

\begin{lstlisting}
#[stack_size(131072)]  // 128KB stack

fn main() -> i32 {
    // Programs with deep recursion need more stack
    return 0
}
\end{lstlisting}

\textbf{Stack Size Guidelines:}

\begin{table}[h]
\centering
\begin{tabular}{ll}
\hline
\textbf{Use Case} & \textbf{Recommended Size} \\
\hline
Simple programs & 8KB -- 16KB \\
Typical applications & 32KB -- 64KB (default) \\
Deep recursion (parsers, tree traversal) & 128KB -- 256KB \\
Large stack arrays & Array size + 32KB buffer \\
\hline
\end{tabular}
\caption{Stack Size Recommendations}
\end{table}

\textbf{Note:} AmigaOS CLI default stack is only 4KB, far too small for most Novus programs. The Novus runtime automatically checks and increases stack size if needed.

\section{Library Attributes}
\label{sec:library-attributes}

These attributes are used when creating AmigaOS shared libraries.

\subsection{@library}

Marks a struct as an Amiga shared library:

\begin{lstlisting}
@library(name = "example.library", version = 1, revision = 0)
pub struct ExampleLibrary {
    base: Library,  // Must include Library base
    counter: u32,
    initialized: bool,
}
\end{lstlisting}

\textbf{Arguments:}
\begin{description}
    \item[\texttt{name}] Library name (must end with \texttt{.library})
    \item[\texttt{version}] Major version number
    \item[\texttt{revision}] Minor revision number
\end{description}

The compiler generates:
\begin{itemize}
    \item ROMTag structure
    \item Library vector table
    \item Open/Close/Expunge handlers
    \item Proper HUNK structure for library loading
\end{itemize}

\subsection{@libfunc}

Marks a method as a library vector entry point:

\begin{lstlisting}
impl ExampleLibrary {
    @libfunc
    pub fn add_numbers(&self, a: i32, b: i32) -> i32 {
        return a + b
    }

    @libfunc
    pub fn get_counter(&self) -> u32 {
        return self.counter
    }
}
\end{lstlisting}

Library functions are called through the library base pointer, following AmigaOS calling conventions.

\subsection{Library Lifecycle Hooks}

\begin{lstlisting}
impl ExampleLibrary {
    @libinit
    pub fn init() -> bool {
        // One-time initialization when library loads
        // Allocate resources, open dependencies
        return true
    }

    @libopen
    pub fn open(version: u32) -> bool {
        // Called each time library is opened
        if version > 1 { return false }  // Version check
        self.counter = self.counter + 1
        return true
    }

    @libclose
    pub fn close() {
        // Called each time library is closed
        self.counter = self.counter - 1
    }

    @libexpunge
    pub fn expunge() -> bool {
        // Return true if library can be unloaded
        // Return false to keep library resident
        return self.counter == 0
    }
}
\end{lstlisting}

\section{Device Attributes}
\label{sec:device-attributes}

These attributes are used when creating AmigaOS devices.

\subsection{@device}

Marks a struct as an Amiga device:

\begin{lstlisting}
@device(name = "example.device", version = 1, revision = 0)
pub struct ExampleDevice {
    base: Device,
    unit: u32,
    active: bool,
}
\end{lstlisting}

\subsection{@devicecmd}

Marks a method as a device command handler:

\begin{lstlisting}
impl ExampleDevice {
    @devicecmd(CMD_READ)
    pub fn handle_read(&self, io: *IORequest) {
        // Handle read command
        io.io_Actual = read_data(io.io_Data, io.io_Length)
    }

    @devicecmd(CMD_WRITE)
    pub fn handle_write(&self, io: *IORequest) {
        // Handle write command
        write_data(io.io_Data, io.io_Length)
        io.io_Actual = io.io_Length
    }
}
\end{lstlisting}

\subsection{@beginio / @abortio}

Standard device I/O handlers:

\begin{lstlisting}
impl ExampleDevice {
    @beginio
    pub fn begin_io(&self, io: *IORequest) {
        // Start I/O operation
        let cmd = io.io_Command
        match cmd {
            CMD_READ => self.handle_read(io),
            CMD_WRITE => self.handle_write(io),
            _ => io.io_Error = IOERR_NOCMD,
        }
    }

    @abortio
    pub fn abort_io(&self, io: *IORequest) -> i32 {
        // Abort pending I/O
        // Return 0 if aborted, error otherwise
        return 0
    }
}
\end{lstlisting}

\subsection{Device Lifecycle Hooks}

Similar to libraries: \texttt{@deviceinit}, \texttt{@deviceopen}, \texttt{@deviceclose}, \texttt{@deviceexpunge}.

\section{Common Attribute Patterns}
\label{sec:common-patterns}

\subsection{Interrupt-Safe State Updates}

Combining \texttt{@atomic} with \texttt{@interrupt\_safe}:

\begin{lstlisting}
static var frame_count: u32 = 0
static var frames_per_second: u32 = 0

@interrupt_safe
@inline
fn increment_frame_count() {
    frame_count = frame_count + 1
}

@atomic
pub fn get_and_reset_fps() -> u32 {
    let fps = frames_per_second
    frames_per_second = 0
    return fps
}

@interrupt
@export
pub fn vblank_interrupt() {
    increment_frame_count()
    if frame_count >= 50 {  // PAL: 50 frames = 1 second
        frames_per_second = frame_count
        frame_count = 0
    }
}
\end{lstlisting}

\subsection{Builder Pattern with Validation}

\begin{lstlisting}
struct ConfigBuilder {
    width: Option<i32>,
    height: Option<i32>,
    depth: i32,
}

impl ConfigBuilder {
    pub fn new() -> ConfigBuilder {
        return ConfigBuilder {
            width: Option::None,
            height: Option::None,
            depth: 8,
        }
    }

    @chain
    pub fn width(&var self, w: i32) {
        self.width = Option::Some(w)
    }

    @chain
    pub fn height(&var self, h: i32) {
        self.height = Option::Some(h)
    }

    @chain
    pub fn depth(&var self, d: i32) {
        self.depth = d
    }

    pub fn build(&self) -> Result<Config, ConfigError> {
        // Validate all required fields are set
        let w = match self.width {
            Option::Some(v) => v,
            Option::None => return Result::Err(ConfigError::MissingWidth),
        }
        let h = match self.height {
            Option::Some(v) => v,
            Option::None => return Result::Err(ConfigError::MissingHeight),
        }
        return Result::Ok(Config { width: w, height: h, depth: self.depth })
    }
}
\end{lstlisting}

\section{Showcase: Interrupt-Driven Music Player}
\label{sec:showcase-music-player}

This showcase demonstrates interrupt and synchronization attributes for audio playback.

\begin{lstlisting}
// music_player_showcase.novus - Demonstrates interrupt attributes

from std::core import Option
from std::io::file import print, println
from amiga::raw::exec import Forbid, Permit, Disable, Enable

// Player state stored as globals for interrupt access
static var g_is_playing: bool = false
static var g_current_position: u32 = 0
static var g_sample_length: u32 = 0
static var g_tempo_counter: u32 = 0
static var g_tempo_divider: u32 = 6  // ~8 ticks/sec at PAL
static var g_volume: u16 = 64

// @interrupt_safe: Documents this is safe to call from interrupt
// @inline: Hint to inline for performance
@interrupt_safe
@inline
fn advance_position() -> bool {
    if g_current_position >= g_sample_length {
        return false
    }
    g_current_position = g_current_position + 1
    return true
}

// @no_interrupts: Wraps in Disable/Enable for hardware access
@no_interrupts
fn write_audio_volume(volume: u16) {
    g_volume = volume
    // In real code: write to Paula register here
}

// @atomic: Wraps in Forbid/Permit for task-safe state access
@atomic
fn start_playback(length: u32, tempo: u32) -> bool {
    if g_is_playing { return false }
    g_sample_length = length
    g_current_position = 0
    g_tempo_divider = tempo
    g_is_playing = true
    return true
}

@atomic
fn stop_playback() {
    if g_is_playing {
        g_is_playing = false
        write_audio_volume(0)
    }
}

// @interrupt: Generates proper save/restore and RTE
// @export: Prevents dead code elimination
@interrupt
@export
pub fn vblank_music_handler() {
    if !g_is_playing { return }

    g_tempo_counter = g_tempo_counter + 1
    if g_tempo_counter < g_tempo_divider { return }
    g_tempo_counter = 0

    if !advance_position() {
        g_is_playing = false
    }
}

pub fn main() -> i32 {
    println("Key Attributes Demonstrated:")
    println("  @interrupt     - VBlank handler")
    println("  @interrupt_safe - Safe helper functions")
    println("  @atomic        - Critical sections")
    println("  @no_interrupts - Hardware access")
    println("  @inline        - Performance hints")
    println("  @export        - Callback preservation")
    return 0
}
\end{lstlisting}

This pattern is typical for:
\begin{itemize}
    \item MOD/XM music players
    \item Sample playback engines
    \item Sound effect systems
    \item Any audio requiring precise timing
\end{itemize}

\section{Attribute Reference}
\label{sec:attribute-reference}

\begin{table}[h]
\centering
\begin{tabular}{llp{6cm}}
\hline
\textbf{Attribute} & \textbf{Applies To} & \textbf{Purpose} \\
\hline
\texttt{@test} & functions & Mark as test case \\
\texttt{@bench} & functions & Mark for benchmarking \\
\texttt{@inline} & functions & Hint to inline \\
\texttt{@noinline} & functions & Prevent inlining \\
\texttt{@packed} & structs & Remove field padding \\
\texttt{@derive} & structs & Auto-implement traits \\
\texttt{@chain} & methods & Enable method chaining \\
\texttt{@export} & functions & Keep and don't mangle name \\
\texttt{@extern\_type} & structs & External C type \\
\texttt{@atomic} & functions & Wrap in Forbid/Permit \\
\texttt{@no\_interrupts} & functions & Wrap in Disable/Enable \\
\texttt{@interrupt} & functions & Mark as interrupt handler \\
\texttt{@interrupt\_safe} & functions & Safe for interrupt context \\
\texttt{@suppress} & functions & Suppress warning \\
\texttt{\#[stack\_size]} & program & Set stack size \\
\texttt{@library} & structs & Define Amiga library \\
\texttt{@libfunc} & methods & Library vector entry \\
\texttt{@device} & structs & Define Amiga device \\
\texttt{@devicecmd} & methods & Device command handler \\
\hline
\end{tabular}
\caption{Complete Attribute Reference}
\end{table}

\section{Unknown Attributes}

If you use an unknown attribute, the compiler emits a warning:

\begin{lstlisting}[style=plain]
warning[W2001]: unknown attribute 'foo'
  --> test.novus:3:1
   |
 3 | @foo
   | ^^^
   |
  help: This attribute is not recognized and will be ignored
  help: Known attributes: library, libfunc, inline, test, ...
\end{lstlisting}

This helps catch typos while allowing forward compatibility with future attributes.

\section{Future Attributes}

The following attributes are planned for future versions:

\begin{description}
    \item[\texttt{@deprecated}] Mark functions as deprecated with migration notes
    \item[\texttt{@must\_use}] Require return value to be used
    \item[\texttt{@cold}] Mark function as unlikely to be called (affects code layout)
    \item[\texttt{@hot}] Mark function as frequently called (affects optimization)
    \item[\texttt{@pure}] Function has no side effects (enables memoization)
\end{description}
